% Entwurf Abschnitt 6.5 (AP7), Stand 20.11.2026. Budget ca. 1,5 Seiten (0,9 / 0,6).
% Code-Stand aller Befunde: G1@25d5c1a. Pfade in % Basis: relativ zu g1pilot/ bzw. Repo-Root.
% Darstellungen (baut ein anderer Auftrag, hier nur referenziert): fig:andockpunkte (Ebenenmodell), tab:andockpunkte.
% N10 (Nils, 16.11.): Bisher ist kein externes Projekt angebunden; 6.5.1 ist rein konzeptionell.
\section{Erweiterungskonzept}
\label{sec:erweiterungskonzept}
Nach Anforderung~A7 soll die Umgebung als Grundlage weiterer Projekte dienen (Tabelle~\ref{tab:anforderungen}).
Abschnitt~\ref{sec:andockpunkte} beschreibt, wo ein externes Projekt andockt, Abschnitt~\ref{sec:container_profile},
wie sich die Umgebung selbst in austauschbare Teile gliedert.

\subsection{Andockpunkte für externe Projekte}
\label{sec:andockpunkte}
% Basis: G1@25d5c1a docs/21-A-Arm-API.md „Überblick“; docs/21-T-Arm-API.md „Sicherheit“,
% „Ohne eigenen HTTP-Client: direkt über DDS“
Externe Projekte unterscheiden sich im Vorwissen und im Bedarf: Eines will nur eine neue Aufgabe einrichten, ein
anderes bringt einen eigenen Regler mit und braucht lediglich Zugang zu Roboter oder Simulation.
Eine einzige Schnittstelle ist für das erste zu anspruchsvoll oder für das zweite zu eng, weshalb die Umgebung
Andockpunkte auf mehreren Ebenen anbietet.
Dabei gilt ein Zielkonflikt: Je tiefer ein Projekt andockt, desto größer ist seine Freiheit und desto weniger
Schutzmechanismen der Umgebung wirken.
Abbildung~\ref{fig:andockpunkte} ordnet die Ebenen nach diesem Zielkonflikt, Tabelle~\ref{tab:andockpunkte} nennt je
Ebene Lieferumfang, Vorwissen und wirksamen Schutz.

% Basis: G1@25d5c1a unitree_mujoco/scene_editor/README.md „Konzept: Basis + Umgebung“, „Hindernis vs. Greif-Objekt“;
% docs/25-T-Eigene-Haende.md „Idee“; docs/30-T-Stehen-und-Gehen.md (Policy-Ordner mit Beschreibungsdatei);
% docs/23-T-Posen-und-Ablaeufe.md
Auf der obersten Ebene erweitert ein Projekt die Umgebung nur über Daten: Szenen, in denen eine Namensregel
Hindernisse von greifbaren Objekten trennt, eigene Hände und weitere Lauf-Policies, die beide nur in der Simulation
wirken, sowie Posen und Abläufe (Abschnitt~\ref{sec:posen_ablaeufe}).
Ohne Programmcode bleiben hier alle Schutzmechanismen wirksam, dafür ist ein Projekt auf die vorhandenen Funktionen
beschränkt.

% Basis: G1@25d5c1a docs/21-A-Arm-API.md „Überblick“; docs/21-T-Arm-API.md „Architektur“, „Sicherheit“, „Parameter“;
% launch/manipulation_launcher.launch.py (arm_api_host 127.0.0.1, arm_api_token leer)
Die zweite Ebene ist eine Programmierschnittstelle der Arme, die Zielposen, Gelenkwinkel und Handstellungen im
JSON-Format über HTTP annimmt, ohne ROS~2 vorauszusetzen, und in der Grundeinstellung nur vom selben Rechner
erreichbar ist.
Sie schreibt auf denselben Kanal wie die Bedienwege und kann die Prüfungen der Armregelung deshalb nicht umgehen; als
eigene Komponente ohne Regelungslogik folgt sie dem Prinzip, Anforderungen in funktional unabhängige Komponenten zu
trennen \cite{Macenski2022}\ev{Macenski2022.04}.

% Basis: G1@25d5c1a docs/02-T-Architektur.md „Zwei Kommunikationsebenen“; docs/31-T-Navigation.md
% „joy_mux / joy_to_cmdvel“, „Warum der Umweg über einen Fake-Joystick?“; docs/12-T-Handy.md „Aktionen und Rechte“
Auf der dritten Ebene nutzt ein Projekt die ROS-2-Kanäle der Umgebung (Abschnitt~\ref{sec:ros2}), über die auch die
Bedienwege senden; so speist die Navigation der Simulation ihre Fahrbefehle über einen Umschalter in denselben Pfad
ein wie die manuelle Steuerung.
Die Schutzmechanismen der Funktionsknoten wirken weiter, Prüfungen einzelner Bedienwege wie Rollen und Steuerhoheit
des Smartphones dagegen nicht, und das Projekt muss ROS~2 sowie Kanäle und Nachrichten der Umgebung kennen.

% Basis: G1@25d5c1a docs/21-T-Arm-API.md „Ohne eigenen HTTP-Client: direkt über DDS“; docs/02-T-Architektur.md
% „Sim ↔ Real: Ein Schalter“
Die unterste Ebene ist die Herstellerschnittstelle.
Mit dem SDK des Herstellers entwickelte Programme laufen auch in unitree\_mujoco \cite{UnitreeMujoco}\ev{UnitreeMujoco.01},
und weil die Simulation der Umgebung dieselbe Schnittstelle bereitstellt (Abschnitt~\ref{sec:sim_real_schalter}),
testet ein vorhandenes SDK-Projekt gegen sie, indem es nur DDS-Domain und Netzwerkschnittstelle umstellt.
Da hier kein Schutzmechanismus der Umgebung wirkt, eignet sich diese Ebene für Projekte mit vollständigem eigenem
Regler einschließlich Kollisionsprüfung.

% Basis: Plan Kap. 6, Abschnitt 5, 6.5.1 (N10)
Ein Projekt dockt daher auf der höchsten Ebene an, die seinen Bedarf deckt.
Ob die konzeptionell angelegten Andockpunkte ein externes Projekt tragen, ist in Abschnitt~\ref{sec:eval_extern}
nachzuweisen, die Aufgabenunabhängigkeit auf der Datenebene in Abschnitt~\ref{sec:eval_aufgaben}.
\pruefen{N10: bisher kein externes Projekt angebunden, Entscheidung zu 8.4.1 offen}
% Umsetzung: sec:entwickler_api, sec:szenen, sec:ersatzregler; sec:hand_profile ergänzen, sobald Kap. 7 es vergibt.
Die Umsetzung beschreiben die Abschnitte~\ref{sec:entwickler_api}, \ref{sec:szenen} und \ref{sec:ersatzregler}.

\subsection{Modularisierung über Container-Profile}
\label{sec:container_profile}
% Basis: G1@25d5c1a docker-compose.yml (Dienste, Profile, gemeinsame Basis-Konfiguration); docs/40-T-Simulation.md „Container“
Nach A6 und A7 soll die Umgebung auf einem neu eingerichteten PC in Betrieb gehen und sich in Teilen austauschen
lassen.
Eine native Installation setzt voraus, dass die Rechenumgebung der ursprünglichen Entwickler nachgebildet werden kann
\cite{Boettiger2015}\ev{Boettiger2015.02}, und eine virtuelle Maschine gilt Kritikern als zu sehr eine Black Box
\cite{Boettiger2015}\ev{Boettiger2015.05} (Abschnitt~\ref{sec:container}).
Die Umgebung nutzt deshalb Container, und zwar statt eines gemeinsamen Images getrennte Container je Aufgabe, die
zwei Profile zusammenstellen.
\pruefen{Begründung gegen ein gemeinsames Image von Nils}

% Basis: G1@25d5c1a docker-compose.yml (profiles sim/real, x-g1pilot-base: network_mode host, Quellstand als Volume);
% docs/40-T-Simulation.md „Container“ (Simulationscontainer ohne ROS); docs/04-T-Umgebungsvariablen.md (Kopf);
% docs/23-T-Posen-und-Ablaeufe.md „Speicherort“; unitree_mujoco/scene_editor/README.md (scenes/)
Der Simulationscontainer kommt ohne ROS~2 aus und nimmt die Stelle des Roboters ein: Das Profil für die Simulation
startet ihn zusammen mit dem Container der Umgebung, das Profil für den Realbetrieb nur die Umgebung.
In beiden Profilen bindet der Container der Umgebung denselben Quellstand ein, und alle Container nutzen das Netzwerk
des Hosts, das die DDS-Kommunikation laut Konfiguration voraussetzt.
Das Verhalten stellen Umgebungsvariablen ein, deren Änderung einen Neustart, aber keinen Neubau der Images erfordert,
und Daten wie Posen, Abläufe und Szenen liegen außerhalb der Container.
Images und Startmenü beschreibt Abschnitt~\ref{sec:bereitstellung}, die Inbetriebnahme bewertet
Abschnitt~\ref{sec:eval_bedienaufwand}.

% Kann
Container lösen nicht jede Frage der Reproduzierbarkeit und des Zeitverhaltens.
Builds eines Dockerfiles sind nicht zwingend bitweise identisch \cite{Boettiger2015}\ev{Boettiger2015.10}, und
Container werden wie native Anwendungen vom Scheduler des Hosts eingeplant \cite{Struhar2020}\ev{Struhar2020.03},
dessen Standard-Linux-Kernel keine zeitlichen Garantien gibt \cite{Struhar2020}\ev{Struhar2020.04}.
